%% kapa: exam=1 round=30 year=2019 period=2 form=A subject=appraisal_law subject_ko=감정평가관계법규 seq=21 number=21
%% source: Q-Net 감정평가사 기출문제 (cst00309) artlSeq=5205205 file=extracted/1008/5205205/2019년 제30회 감정평가사 1차시험 문제지/2019년 제30회 감정평가사 1차시험 문제지/2019년도 제30회 감정평가사 1차 2교시 A형.hwp
%% conversion: kordoc 4.24.51 -> tools/gen1.py (auto); review=proofread-sonnet; flags=none
\begin{question}{21}
국유재산법령상 무단점유자에 관한 설명으로 옳지 않은 것은?
\begin{choices}
  \item 행정재산에 대한 사용허가 기간이 끝난 후 다시 사용허가 없이 행정재산을 계속 사용한 자는 무단점유자에 해당한다.
  \item 정당한 사유 없이 국유재산에 시설물을 설치한 경우 중앙관서의 장등은「행정대집행법」을 준용하여 철거할 수 있다.
  \item 무단점유자가 재해로 재산에 심한 손실을 입은 경우는 중앙관서의 장등이 변상금 징수를 미룰 수 있는 사유에 해당한다.
  \item 변상금의 연체료 부과대상이 되는 연체기간은 납기일부터 60개월을 초과할 수 없다.
  \item 중앙관서의 장등은 행정재산의 무단점유자에 대하여 그 재산에 대한 사용료의 100분의 150에 상당하는 변상금을 징수한다.
\end{choices}
\end{question}
\answer{5}
